\chapter{Marco referencial y antecedentes}
\label{cap:marco-referencial}

\section{Antecedentes del proyecto}
\label{sec:antecedentes}

El presente trabajo de grado se desarrolla en la Universidad Antonio Nariño (UAN) como continuación del anteproyecto titulado \emph{Protocolo distribuido de control y reconfiguración topológica con continuidad lógica para redes multisalto sobre Wi-Fi convencional}. El anteproyecto estableció el planteamiento, objetivos preliminares y alcance; este documento profundiza en diseño formal, implementación y validación experimental.

Los autores ---Juan Carlos Clavijo Triviño y Brandon Stiven Ganzo Murcia--- dividieron responsabilidades según la arquitectura del sistema: conectividad Wi-Fi (GO, DNS-SD, STA, plano TCP) y enrutamiento/topología (tablas, heartbeats, reenvío).

\section{Prototipo EXP-01 (proto-exp01-gostat)}
\label{sec:exp01}

Como antecedente de implementación se desarrolló el prototipo experimental EXP-01, documentado en \texttt{docs/prototype/proto-exp01-gostat/}. Este prototipo demostró en laboratorio:

\begin{itemize}
  \item Formación de grupo P2P GO en dos dispositivos Android independientes.
  \item Descubrimiento DNS-SD del servicio \texttt{\_pct-ctrl} con parseo de TXT.
  \item Asociación STA legacy del hijo al SoftAP del padre.
  \item Coexistencia GO propio + STA upstream (hipótesis H1).
\end{itemize}

EXP-01 identificó restricciones operativas críticas: el buscador no debe mantener GO activo durante el escaneo DNS-SD; y la identidad de hijos no debe depender de \texttt{clientList} MAC cuando exista handshake formal.

\section{Evolución hacia lib-pct-core}
\label{sec:evolucion-core}

El prototipo EXP-01 se refactorizó en la biblioteca Maven \texttt{co.uan.pct:core} (\texttt{pct/library/lib-pct-core}), con API pública \texttt{PctNode}:

\begin{lstlisting}[language=Java,caption={API pública PCT (Kotlin)},label={lst:api-pct}]
val pct = PctCore.create()
pct.init(context, PctConfig.DEFAULT)
pct.start()  // bootstrap: scan -> join | root
pct.topology.collect { ... }
pct.sendUser(destinationNid, payload)
pct.close()
\end{lstlisting}

La biblioteca concentra bootstrap automático, teardown P2P anti-zombi al cerrar la aplicación, y observables (\texttt{StateFlow}) de fase, topología, vecinos y rutas.

\section{Requisitos funcionales}
\label{sec:requisitos-funcionales}

La Tabla~\ref{tab:rf} resume los requisitos funcionales trazables (RF). Los marcados \emph{Must} son obligatorios para la versión de grado.

\begin{longtable}{@{}clp{7cm}@{}}
  \caption{Requisitos funcionales PCT}
  \label{tab:rf} \\
  \toprule
  \textbf{ID} & \textbf{Pri.} & \textbf{Descripción} \\
  \midrule
  \endfirsthead
  \toprule
  \textbf{ID} & \textbf{Pri.} & \textbf{Descripción} \\
  \midrule
  \endhead
  RF-01 & Must & Identidad UUID única por nodo \\
  RF-02 & Must & Arranque homogéneo como P2P GO \\
  RF-03 & Must & Prohibición de \texttt{connect()} para topología \\
  RF-04 & Must & Enlace upstream STA legacy al padre \\
  RF-05 & Must & Anuncio DNS-SD en GO \\
  RF-06 & Must & Unión simétrica (\texttt{\_pct-seek} / \texttt{JOIN\_OFFER}) \\
  RF-07 & Must & Handshake TCP post-asociación (HELLO + JOIN\_COMMIT) \\
  RF-08 & Must & Tabla de rutas local por \texttt{node\_id} \\
  RF-09 & Must & Prevención de ciclos (\texttt{path\_trace}, \texttt{hop\_limit}) \\
  RF-10 & Must & Heartbeats PING/PONG \\
  RF-11 & Must & Detección de fallo y NODE\_DOWN \\
  RF-12 & Must & Reconfiguración sin reinicio global \\
  RF-13 & Should & Multisalto $\geq$ 3 nodos \\
  RF-14 & Should & Unión tardía sin reiniciar red existente \\
  RF-15 & Should & Continuidad \texttt{session\_epoch} \\
  RF-16 & Could & DAG lógico multipath \\
  RF-17 & Won't & Cifrado E2E (v1) \\
  RF-18 & Won't & QoS (v1) \\
  \bottomrule
\end{longtable}

\section{Requisitos no funcionales}
\label{sec:rnf}

\begin{itemize}
  \item[RNF-01] Plataforma mínima Android~12 (API~31).
  \item[RNF-02] Wi-Fi Direct mandatorio.
  \item[RNF-03] RAM mínima 2~GB.
  \item[RNF-04] Profundidad máxima 7 saltos.
  \item[RNF-05] Reconexión objetivo $\leq$ 15~s.
  \item[RNF-06] Puerto control TCP 8765 (configurable).
  \item[RNF-07] Codificación wire binaria big-endian.
  \item[RNF-08--09] Payload TCP $\leq$ 1400~B; mensaje usuario $\leq$ 1024~B.
\end{itemize}

\section{Restricciones de plataforma Android}
\label{sec:restricciones-android}

Android~12+ impone permisos \texttt{NEARBY\_WIFI\_DEVICES} (API~33+) y restricciones de descubrimiento P2P concurrente con GO activo. No es posible modificar tramas 802.11 ni implementar NAT multisalto en kernel sin \emph{root}. PCT opera íntegramente en espacio de usuario mediante \texttt{Network.bindSocket()} para seleccionar la interfaz STA upstream frente al GO propio.

\section{Perfiles de hardware}
\label{sec:perfiles-hardware}

Los dispositivos de laboratorio se clasifican en perfiles $\alpha$, $\beta$, $\gamma$ según capacidad dual GO+STA concurrente, RAM y batería. El perfil mínimo $\alpha$ exige verificación de hipótesis H1 (GO y STA simultáneos) antes de continuar experimentos multisalto.

\section{Síntesis del capítulo}
\label{sec:sintesis-marco-referencial}

El marco referencial vincula el anteproyecto, el prototipo EXP-01 y la biblioteca \texttt{lib-pct-core} con un conjunto trazable de requisitos y restricciones Android que condicionan todas las decisiones de diseño documentadas en el Capítulo~\ref{cap:analisis-diseno}.
